\documentclass[preview]{standalone}

\usepackage{amsmath}
\usepackage{amssymb}
\usepackage{stellar}
\usepackage{definitions}

\begin{document}

\id{coprime-ideals-chinese-remainder}
\genpage

\section{Comaximal ideals}

\begin{snippetdefinition}{comaximal-ideals-definition}{Comaximal ideals}
    Let \(A\) be a \commutativering and let \(I,J\idealin A\). The
    \ideal[ideals] \(I\) and \(J\) are \emph{comaximal} if
    \[
        I+J=A.
    \]
    Equivalently, there exist \(x\in I\) and \(y\in J\) such that
    \[
        x+y=1_A.
    \]
\end{snippetdefinition}

\begin{snippetexample}{integer-comaximal-ideals-example}{Comaximal ideals in the integers}
    If \(m,n\in\integers\) satisfy \(\gcd(m,n)=1\), then the
    \ideal[ideals] \((m)\) and \((n)\) are comaximal. Indeed, Bézout's
    identity gives \(h,k\in\integers\) such that
    \[
        hm+kn=1,
    \]
    and therefore \(1\in(m)+(n)\).
\end{snippetexample}

\begin{snippetproposition}{powers-of-comaximal-ideals}{Powers of comaximal ideals are comaximal}
    Let \(A\) be a \commutativering and let \(I,J\idealin A\) be
    \comaximal. For every \(n,m\geq1\),
    \[
        I^n+J^m=A.
    \]
\end{snippetproposition}

\begin{snippetproof}{powers-of-comaximal-ideals-proof}{powers-of-comaximal-ideals}{Powers of comaximal ideals are comaximal}
    Choose \(x\in I\) and \(y\in J\) such that \(x+y=1_A\), and put
    \(s=n+m-1\). In the binomial expansion
    \[
        1_A=(x+y)^s
        =
        \sum_{h=0}^{s}\binom{s}{h}x^hy^{s-h},
    \]
    every term belongs to \(I^n\) or to \(J^m\). Indeed, if \(h\geq n\)
    it belongs to \(I^n\), whereas if \(h<n\), then
    \(s-h\geq m\), so it belongs to \(J^m\). Hence
    \(1_A\in I^n+J^m\).
\end{snippetproof}

\begin{snippetproposition}{ideal-product-intersection-comaximal}{Product and intersection of pairwise comaximal ideals}
    Let \(A\) be a \commutativering, and let
    \(I_1,\ldots,I_r\idealin A\) be pairwise \comaximal. Then
    \[
        I_1\cdots I_r
        =
        \bigcap_{i=1}^{r}I_i.
    \]
\end{snippetproposition}

\begin{snippetproof}{ideal-product-intersection-comaximal-proof}{ideal-product-intersection-comaximal}{Product and intersection of pairwise comaximal ideals}
    First let \(r=2\). One always has
    \(I_1I_2\subseteq I_1\intersection I_2\). Choose
    \(x\in I_1\) and \(y\in I_2\) with \(x+y=1_A\). If
    \(z\in I_1\intersection I_2\), then
    \[
        z=z(x+y)=zx+zy\in I_1I_2,
    \]
    proving \(I_1I_2=I_1\intersection I_2\).

    For the general case, proceed by induction. Pairwise \comaximal[comaximality]
    implies
    \[
        I_r+\bigcap_{i=1}^{r-1}I_i=A.
    \]
    To see this, choose \(x_i\in I_r\) and \(y_i\in I_i\) with
    \(x_i+y_i=1_A\), and expand
    \[
        1_A=\prod_{i=1}^{r-1}(x_i+y_i).
    \]
    Every term except \(y_1\cdots y_{r-1}\) lies in \(I_r\), while that
    exceptional term lies in every \(I_i\), for \(i<r\). Applying the
    two-ideal case and the induction hypothesis yields
    \begin{align*}
        I_1\cdots I_r
        &=
        I_r\left(\bigcap_{i=1}^{r-1}I_i\right)\\
        &=
        I_r\intersection\bigcap_{i=1}^{r-1}I_i
        =
        \bigcap_{i=1}^{r}I_i.
    \end{align*}
\end{snippetproof}

\section{Chinese remainder theorem}

\begin{snippetexample}{finite-direct-product-rings-example}{Finite direct products of rings}
    Let \(A_1,\ldots,A_r\) be \commutativering[commutative rings]. Their
    direct product
    \[
        \prod_{i=1}^{r}A_i
    \]
    is a \commutativering under componentwise addition and multiplication.
    In particular, for two factors,
    \[
        (a_1,b_1)+(a_2,b_2)
        =
        (a_1+a_2,b_1+b_2)
    \]
    and
    \[
        (a_1,b_1)(a_2,b_2)
        =
        (a_1a_2,b_1b_2).
    \]
    Its zero and identity are
    \[
        (0_{A_1},\ldots,0_{A_r}),
        \qquad
        (1_{A_1},\ldots,1_{A_r}).
    \]
\end{snippetexample}

\begin{snippettheorem}{commutative-ring-chinese-remainder-theorem}{Chinese remainder theorem for commutative rings}
    Let \(A\) be a \commutativering and let
    \(I_1,\ldots,I_r\idealin A\) be pairwise \comaximal. The
    \ringhomomorphism
    \[
        \varphi\colon A\fromto\prod_{i=1}^{r}A/I_i,
        \qquad
        \varphi(a)=(a+I_1,\ldots,a+I_r),
    \]
    is \surjective and satisfies
    \[
        \ringker\varphi
        =
        \bigcap_{i=1}^{r}I_i
        =
        I_1\cdots I_r.
    \]
    Consequently,
    \[
        A/(I_1\cdots I_r)
        \ringisomorphic
        \prod_{i=1}^{r}A/I_i.
    \]
\end{snippettheorem}

\begin{snippetproof}{commutative-ring-chinese-remainder-theorem-proof}{commutative-ring-chinese-remainder-theorem}{Chinese remainder theorem for commutative rings}
    Since each component of \(\varphi\) is a canonical projection,
    \(\varphi\) is a \ringhomomorphism. Moreover,
    \[
        \ringker\varphi
        =
        \{a\in A\suchthat a\in I_i\text{ for all }i\}
        =
        \bigcap_{i=1}^{r}I_i
        =
        I_1\cdots I_r.
    \]

    For every \(i\), the \ideal \(I_i\) is \comaximal with
    \(\bigcap_{j\neq i}I_j\). Choose
    \[
        x_i\in I_i,
        \qquad
        y_i\in\bigcap_{j\neq i}I_j,
        \qquad
        x_i+y_i=1_A.
    \]
    Then \(y_i\equiv1_A\pmod{I_i}\) and
    \(y_i\equiv0_A\pmod{I_j}\) for \(j\neq i\). Given
    \[
        (b_1+I_1,\ldots,b_r+I_r)
        \in\prod_{i=1}^{r}A/I_i,
    \]
    the element \(b=\sum_{i=1}^{r}b_iy_i\) satisfies
    \[
        \varphi(b)=(b_1+I_1,\ldots,b_r+I_r).
    \]
    Thus \(\varphi\) is \surjective. The first isomorphism theorem for
    rings gives the asserted isomorphism.
\end{snippetproof}

\begin{snippetcorollary}{commutative-ring-crt-congruences-corollary}{Chinese remainder theorem for congruences}
    Let \(n_1,\ldots,n_r\in\integers\) be pairwise \coprime positive
    integers. For every \(b_1,\ldots,b_r\in\integers\), the system
    \[
        x\equiv b_i\pmod{n_i},
        \qquad
        i=1,\ldots,r,
    \]
    has a solution, unique modulo \(N=n_1\cdots n_r\).
\end{snippetcorollary}

\begin{snippetproof}{commutative-ring-crt-congruences-corollary-proof}{commutative-ring-crt-congruences-corollary}{Chinese remainder theorem for congruences}
    The \ideal[ideals] \((n_i)\) are pairwise \comaximal by Bézout's
    identity. Hence the ring-theoretic Chinese remainder theorem gives a
    \surjective \ringhomomorphism
    \[
        \integers\fromto
        \prod_{i=1}^{r}\integers/n_i\integers.
    \]
    Surjectivity gives existence of a simultaneous solution, while its
    kernel is
    \[
        \bigcap_{i=1}^{r}(n_i)
        =
        (n_1\cdots n_r).
    \]
    Thus two solutions are congruent modulo \(N\).
\end{snippetproof}

\end{document}
